% Lösungen Einheit 3 von 4 – Ebenen im Koordinatensystem (zusammenbau v0.1)
\einheitenkopf{Einheit 3 von 4 · Ebenen im Koordinatensystem}

\erg{15}{a) $y$ \quad b) $y = 0$ \quad c) $x$ kommt nicht vor: $\vec n = (0 | 2 | 3)$ steht senkrecht auf der Richtung der $x$-Achse, $E$ ist also parallel zu ihr; sie enthält die Achse nicht, weil der Ursprung $0 \ne 18$ liefert \quad d) $z$ fehlt und der Ursprung erfüllt die Gleichung: $E$ enthält die $z$-Achse, ist also nicht nur parallel zu ihr \quad e) $S_x(2 | 0 | 0)$, $S_y(0 | 6 | 0)$, $S_z(0 | 0 | 3)$ \quad f) die $y$-Koordinate ist immer 4: $L\colon y = 4$, parallel zur $xz$-Ebene (nicht die $xz$-Ebene selbst) \quad g) nein: die $yz$-Ebene hat die Gleichung $x = 0$; $L$ enthält zwar den Ursprung, ist aber eine andere Ebene – $(5 | 6 | 0)$ liegt in $L$, nicht in der $yz$-Ebene \quad h) Der Richtungsvektor von $g$ ist Normalenvektor der Rechtecksebene und zeigt in $y$-Richtung, die Ebene hat also die Form $y = c$; aus $P$ folgt $c = 0$: die $xz$-Ebene \quad i) $\overrightarrow{AB} = (4 | -2 | 0)$, $\overrightarrow{AD} = (-1 | 0{,}5 | 4)$, $\vec n = (1 | 2 | 0)$: die dritte Komponente ist 0, der Normalenvektor liegt waagerecht, die Wand steht senkrecht auf dem Boden \quad j) Der Normalenvektor der $xz$-Ebene zeigt in $y$-Richtung; eine Ebene, die senkrecht auf der $xz$-Ebene steht, hat einen Normalenvektor, der dazu orthogonal ist – seine $y$-Komponente ist 0, also $k = 0$}

15e)

\begin{ksys3}[xyz]\rebene{2}{6}{3}{E}\end{ksys3}

\erg{16}{a) Der Ursprung erfüllt die Gleichung nicht ($0 \ne 30$); $E$ ist nur parallel zur $y$-Achse \quad b) Erfüllt ein Punkt die Gleichung, dann auch jeder Punkt, der sich nur in der $z$-Koordinate von ihm unterscheidet; mit jedem Punkt liegt also die ganze Parallele zur $z$-Achse in der Ebene. \quad c) ja, $z$ fehlt, $\vec n = (4 | 3 | 0)$ liegt waagerecht; die Unterkante trifft die $x$-Achse in $(9 | 0 | 0)$ \quad d) ein Rechteck in der Höhe 3 parallel zur $xy$-Ebene, z. B. mit den Ecken $(0 | 0 | 3)$, $(4 | 0 | 3)$, $(4 | 5 | 3)$, $(0 | 5 | 3)$}

16d)

\begin{ksys3}[xyz]\rebenepar*{0,0,3}{4,0,0}{0,5,0}{E}\end{ksys3}

